\subsection{کالبدشکافی ماژول \lr{rala\_compressor}}

\subsubsection*{۱. فلسفه طراحی، پارادایم معماری و تمثیل ملموس}
ماژول فشرده‌ساز \lr{rala\_compressor} فرآیند تبدیل وزن‌های حجیم مدل‌های زبانی را به ساختار بازگشتی رتبه‌افزوده محقق می‌سازد. این موتور با استخراج الگوهای اصلی و جبران خطای باقیمانده، وزن‌ها را برای استنتاج فوق‌سبک آماده می‌کند.

\textbf{تمثیل ملموس و مثال کاربردی:}
تبدیل یک تابلوی نقاشی رنگ روغن با ابعاد ۳ متر در ۲ متر به یک پوستر دیجیتال باکیفیت: ابتدا بافت عمومی و طیف رنگی پس‌زمینه با یک اسکن کلی ملایم ثبت می‌شود (تقریب خطی)، سپس لبه‌های پرتره و ریزه‌کاری‌های چهره با لایه‌های برداری وکتور با رزولوشن میکرومتری اضافه می‌گردند (افزونه کم‌رتبه).

\subsubsection*{۲. فرمول‌بندی و استنتاج عمیق ریاضی \lr{(Mathematical Derivation)}}
تانسورهای توجه به فرم بازگشتی تبدیل شده و عملگر خطای تقریب با تجزیه مقادیر تکین مانده بهینه می‌شود:
\begin{equation}
\mathbf{E} = \mathbf{W}_{\text{attn}} - \mathbf{W}_{\text{linear}}, \quad \mathbf{E} \approx \sum_{k=1}^r \mathbf{u}_k \mathbf{v}_k^\top
\end{equation}
با بهینه‌سازی توأم پارامترها در فاز تنظیم دقیق، خطای بازسازی به حداقل می‌رسد:
\begin{equation}
\mathcal{L}_{\text{distill}} = \|\mathbf{Y}_{\text{orig}} - (\mathbf{Y}_{\text{linear}} + \mathbf{U}_r \mathbf{V}_r^\top \mathbf{X})\|_2^2
\end{equation}
بدین شکل وزن‌های متراکم بدون افت حافظه جایگزین می‌گردند.

\subsubsection*{۳. کالبدشکافی الگوریتمی و جریان داده \lr{(Algorithmic Tensor Trace)}}
در پیاده‌سازی سورس پایتون در فایل ماژول \lr{rala\_compressor.py}، پردازش در متد \lr{forward} به صورت مرحله‌به‌مرحله و گام‌به‌گام صورت می‌پذیرد:
\begin{enumerate}
\item واکشی ماتریس وزن‌های پیش‌آموزش‌دیده از فایل چک‌پوینت مدل.
\item فاکتورسازی پایه‌ای و استخراج ماتریس‌های تبدیل خطی اولیه.
\item محاسبه ماتریس باقیمانده خطا و استخراج $r=16$ بردار ویژه اصلی.
\item کالیبراسیون وزن‌ها بر روی داده‌های متنی با الگوریتم تقطیر دانش.
\item بسته‌بندی وزن‌ها در قالب ماژول فشرده آماده استنتاج.
\end{enumerate}

\subsubsection*{۴. تحلیل همجوشی سخت‌افزاری و ترافیک حافظه \lr{(Hardware \& Roofline Analysis)}}
از دیدگاه مدل سقف کارایی (\lr{Roofline Model}) و معماری پردازنده‌های گرافیکی مدرن (مانند \lr{NVIDIA Ada / RTX 4070}):
\begin{itemize}
\item \textbf{انقباض فیزیکی ردپای وزن‌ها:} اندازه حافظه دائمی مورد نیاز برای ذخیره وزن‌ها به میزان ۶۵ درصد کاهش می‌یابد.
\item \textbf{حذف بارگذاری مکرر از حافظه اصلی:} کلیه محاسبات لایه فشرده در کش پایدار پردازنده انجام می‌شود.
\end{itemize}

\subsubsection*{۵. ارزیابی تجربی، بنچمارک و توجیه علمی نتایج}
بر اساس آزمون‌های ثبت‌شده در خط لوله ارزیابی پایان‌نامه:
\begin{itemize}
\item \textbf{وضعیت رسمی آزمون کیفی:} \lr{PASS}
\item \textbf{دقت ثبت‌شده روی بنچ‌مارک:} 60.09 درصد
\item \textbf{شاخص ارزیابی آماری (\lr{F1 Score}):} 62.3 درصد
\item \textbf{اوج مصرف حافظه گرافیکی (\lr{Peak VRAM}):} 750.0 مگابایت (در برابر ۴۲۰۰ مگابایت مدل مبنا)
\item \textbf{نرخ توان عملیاتی (\lr{Throughput}):} 36,500 توکن بر ثانیه
\end{itemize}

\textbf{توجیه علمی و علت‌شناسی جذاب نتایج:}
ثبت مصرف VRAM پایدار در سطح 750.0 مگابایت (کاهش 82.1 درصدی در مقایسه با ۴۲۰۰ مگابایت مبنا) و توان پردازش 36,500 توکن در ثانیه، نمایانگر کارایی این ماژول در فشرده‌سازی عمیق وزن‌ها برای محیط‌های با حافظه محدود است.

\subsubsection*{۶. دیاگرام معماری شکیل در لاتک \lr{(TikZ Architecture Flowchart)}}
جهت تبیین دقیق جریان داده تانسوری و نحوه تغییر شکل ماتریس‌ها، نمودار معماری این ماژول با استفاده از بسته قدرتمند \lr{TikZ} در محیط لاتک ترسیم شده است:

\begin{center}
\begin{tikzpicture}[scale=0.82, every node/.style={scale=0.82}, 
    node distance=1.1cm,
    auto,
    block/.style={rectangle, draw=blue!70, fill=blue!12, rounded corners=4pt, text centered, minimum height=0.85cm, font=\small},
    comp/.style={rectangle, draw=orange!80, fill=orange!18, rounded corners=4pt, text centered, minimum height=0.85cm, font=\small},
    outblock/.style={rectangle, draw=teal!80, fill=teal!15, rounded corners=4pt, text centered, minimum height=0.85cm, font=\small}
]
    \node [block] (input) {وزن مبنا $\mathbf{W}$};
    \node [comp, right=1.0cm of input] (lin) {پایه خطی};
    \node [comp, right=1.0cm of lin] (res) {مانده کم‌رتبه};
    \node [outblock, right=1.0cm of res] (out) {وزن فشرده \lr{RALA}};

    \draw [->, thick, draw=blue!80] (input) -- (lin);
    \draw [->, thick, draw=orange!80] (lin) -- (res);
    \draw [->, thick, draw=orange!80] (res) -- (out);
\end{tikzpicture}
\end{center}

\subsubsection*{۷. پیوست تصویری و پرامپت تولید هوش مصنوعی}
برای ترسیم تجسم سه‌بعدی و گرافیک صنعتی این مکانیزم با ابزارهای هوش مصنوعی، پرامپت انگلیسی زیر با دقت تدوین شده است:

\vspace{0.2cm}
\begin{latin}
\footnotesize
\textbf{AI Technical Visualization Prompt:}\\
"A 3D technological visualization of the RALA compressor. An intricate crystalline neural network layer being systematically compressed into a streamlined dual-core processor core, cyan data pulses, 8k resolution --ar 16:9"
\end{latin}

\begin{figure}[H]
    \centering
    \framebox[0.9\textwidth]{
    \begin{minipage}{0.85\textwidth}
    \vspace{1.8cm}
    \centering
    \textbf{[کادر الحاق تصویر گرافیکی سه‌بعدی ماژول \lr{rala\_compressor}]}\\[0.2cm]
    \small (بر اساس پرامپت فنی فوق تولید و در این کادر جایگزین می‌شود)
    \vspace{1.8cm}
    \end{minipage}
    }
    \caption{شماتیک و جریان داده تانسوری در ماژول \lr{rala\_compressor}}
    \label{fig:rala_compressor}
\end{figure}

\vspace{0.5cm}
\hrule
\vspace{0.5cm}
